% ==============================================================================
% Chapter 11: Single-Cell RNA Sequencing
% ==============================================================================
\chapter{Single-Cell RNA Sequencing}
\label{ch:scrnaseq}

\begin{keyconcept}
Single-cell RNA sequencing (\acrshort{scRNA-seq}) measures the transcriptome
of individual cells, revealing cellular heterogeneity that is invisible to bulk
RNA-seq. It has transformed our understanding of development, immunology,
neuroscience, and cancer by enabling the identification of rare cell types,
continuous differentiation trajectories, and cell-state-specific gene
regulatory programs.
\end{keyconcept}

% ---------------------------------------------------------------------------
\section{Why Single-Cell Resolution?}
\label{sec:scrnaseq:why}

\subsection{The Limitation of Bulk RNA-seq}

Bulk RNA-seq measures the \emph{average} gene expression across a population
of $10^5$--$10^7$ cells. This averaging has profound consequences:

\begin{itemize}[nosep]
  \item \textbf{Masking of rare populations:} A cell type comprising 1\% of a
    tissue may be functionally critical (e.g., stem cells, tumor-initiating
    cells) but its transcriptional signature is diluted beyond detection in
    bulk data.
  \item \textbf{False intermediate states:} A mixture of two distinct cell
    states (e.g., activated and resting T cells) appears in bulk as a single
    ``intermediate'' state that does not actually exist in any individual cell.
  \item \textbf{Ecological fallacy:} Correlations observed at the bulk level
    may not hold at the single-cell level. Two genes may appear co-expressed
    in bulk data simply because they are expressed in different subpopulations
    that co-vary in proportion.
  \item \textbf{No trajectory information:} Bulk RNA-seq provides a snapshot
    but cannot reveal the order of events during dynamic processes such as
    differentiation, reprogramming, or drug response.
\end{itemize}

\subsection{What Single-Cell Resolution Enables}

\begin{enumerate}[nosep]
  \item \textbf{Cell type discovery:} Unbiased identification of cell types
    and subtypes without prior marker knowledge.
  \item \textbf{Developmental trajectories:} Reconstruction of differentiation
    paths and branching decisions using pseudotime analysis.
  \item \textbf{Gene regulatory networks:} Inference of transcription factor
    activity and regulatory relationships within specific cell types.
  \item \textbf{Disease heterogeneity:} Characterization of tumor
    subclones, immune microenvironment composition, and treatment-resistant
    cell populations.
  \item \textbf{Cell atlases:} Comprehensive catalogues of all cell types in
    an organism (Human Cell Atlas, Tabula Sapiens).
\end{enumerate}


% ---------------------------------------------------------------------------
\section{A Brief History of scRNA-seq}
\label{sec:scrnaseq:history}

\begin{description}[style=nextline]
  \item[2009 --- Tang et al.] First scRNA-seq experiment. Manual isolation of
    single mouse blastomeres, poly-A-tailed cDNA amplification, sequencing on
    the SOLiD platform. Detected $\sim$5,000 genes per cell.
  \item[2011 --- Islam et al.] Introduced STRT-seq with template-switching for
    5$'$ end capture and early multiplexing using cell barcodes.
  \item[2012 --- Hashimoto et al. / Ramsk\"{o}ld et al.] Smart-seq protocol: full-length
    transcript coverage from single cells, greatly improving gene body coverage
    and isoform detection.
  \item[2014 --- Smart-seq2 (Picelli et al.)] Improved sensitivity and reduced
    cost compared to Smart-seq. Became the gold standard for plate-based
    scRNA-seq.
  \item[2015 --- Drop-seq \& inDrop] Droplet-based microfluidics enabled
    profiling of thousands of cells in a single experiment. Drop-seq (Macosko
    et al.) used barcoded beads in nanoliter droplets. inDrop (Klein et al.)
    used hydrogel beads. Both introduced the concept of \acrshort{UMI}-based
    counting.
  \item[2016 --- 10x Genomics Chromium] Commercialized droplet-based scRNA-seq
    with the GemCode/Chromium platform. Standardized reagents, automated
    workflow, and the \software{Cell Ranger} analysis pipeline made scRNA-seq
    accessible to any lab with a sequencer.
  \item[2017--2019 --- Scale-up era] sci-RNA-seq (combinatorial indexing),
    SPLiT-seq, SHARE-seq, and Slide-seq extended throughput to millions of
    cells and added spatial and multi-omic dimensions.
  \item[2020--present --- Atlas and multi-omic era] The Human Cell Atlas
    project, COVID-19 single-cell studies, and multi-omic platforms (10x
    Multiome, CITE-seq, ASAP-seq) have made scRNA-seq a routine component of
    modern biology.
\end{description}


% ---------------------------------------------------------------------------
\section{Cell Isolation Methods}
\label{sec:scrnaseq:isolation}

The first step in any scRNA-seq experiment is isolating individual cells and
capturing their RNA. The choice of isolation method determines throughput,
sensitivity, and cost.

\subsection{FACS Sorting}

Fluorescence-activated cell sorting (FACS) uses fluorescent antibodies to sort
cells one-by-one into wells of a 96- or 384-well plate. Each well contains
lysis buffer and can be processed with a plate-based scRNA-seq protocol
(Smart-seq2, Smart-seq3).

\textbf{Advantages:} Protein-level selection of specific populations (e.g.,
CD8$^+$ T cells); index sorting records surface marker levels for each cell.
\textbf{Limitations:} Low throughput (hundreds to low thousands of cells);
requires viable single-cell suspensions and surface markers.

\subsection{Microfluidics}

The Fluidigm C1 system captures cells in microfluidic chambers using
integrated fluidic circuits. Each chip captures 96 or 800 cells.

\textbf{Advantages:} Automated cell capture and lysis; full-length protocols
(Smart-seq).
\textbf{Limitations:} Low throughput; size-dependent capture (requires cells of
uniform size); high cost per cell; largely superseded by droplet-based methods.

\subsection{Droplet-Based Methods}
\label{sec:scrnaseq:droplet}

Droplet microfluidics encapsulates individual cells with barcoded beads in
nanoliter oil droplets, enabling massively parallel scRNA-seq.

\begin{description}[style=nextline]
  \item[10x Genomics Chromium] The dominant commercial platform (see
    \cref{sec:scrnaseq:10x} for details). Uses Gel Bead-in-Emulsions (GEMs).
    Throughput: 500--10,000+ cells per reaction.
  \item[Drop-seq] Academic protocol using barcoded microparticles and a simple
    PDMS microfluidic device. Lower cost per cell than 10x but higher doublet
    rate and lower capture efficiency.
  \item[inDrop] Uses deformable hydrogel beads for more uniform barcode
    delivery. Open-source design.
\end{description}

\subsection{Combinatorial Indexing}

These methods achieve single-cell resolution without physically isolating
individual cells. Instead, cells are split across wells, indexed with
well-specific barcodes in multiple rounds, and then pooled. The extremely
low probability that two cells receive the same combination of barcodes across
all rounds provides unique identification.

\begin{description}[style=nextline]
  \item[sci-RNA-seq / sci-RNA-seq3] Achieves profiling of millions of cells in
    a single experiment. Uses two or three rounds of combinatorial indexing
    (reverse transcription barcoding $\rightarrow$ ligation barcoding
    $\rightarrow$ PCR barcoding).
  \item[SPLiT-seq] Split-pool ligation-based transcriptome sequencing. Uses
    successive rounds of in-cell barcoding by ligation. No specialized
    equipment required beyond standard liquid handling.
\end{description}

\begin{tipbox}
Combinatorial indexing methods excel for very large experiments (e.g.,
organism-scale cell atlases). They are cost-effective per cell but have lower
sensitivity (fewer genes per cell) compared to droplet-based or plate-based
methods.
\end{tipbox}

\subsection{Plate-Based Methods}

\begin{description}[style=nextline]
  \item[Smart-seq2] Full-length transcript coverage from single cells sorted
    into wells. Uses oligo-dT priming, template switching, and limited PCR
    amplification. Highest sensitivity ($\sim$5,000--10,000 genes per cell) and
    gene body coverage of any scRNA-seq method. Enables isoform-level analysis
    and SNV detection. Throughput: 96--384 cells per run.
  \item[Smart-seq3 / Smart-seq3xpress] Adds UMIs to Smart-seq2 while retaining
    full-length information. Combines the quantitative accuracy of UMI counting
    with the isoform resolution of full-length coverage.
\end{description}


% ---------------------------------------------------------------------------
\section{10x Genomics Chromium in Detail}
\label{sec:scrnaseq:10x}

The 10x Genomics Chromium platform is the most widely used scRNA-seq system
worldwide. Understanding its mechanics is essential for interpreting scRNA-seq
data.

% TikZ diagram: 10x Chromium droplet-based workflow
\begin{figure}[htbp]
\centering
\begin{tikzpicture}[
  node distance=0.6cm and 0.3cm,
  every node/.style={font=\small},
  boxstep/.style={draw, rounded corners, minimum width=2.8cm, minimum height=1.0cm, align=center, font=\small},
  gem/.style={draw, circle, minimum size=1.8cm, fill=yellow!15, font=\scriptsize, align=center},
  arrow/.style={-{Stealth[length=5pt]}, thick, MidnightBlue!70},
  label/.style={font=\scriptsize\itshape, text=gray!60!black},
]

% Step 1: Cell suspension + Gel Beads + Oil
\node[boxstep, fill=blue!10] (cells) {Cell\\suspension};
\node[boxstep, fill=orange!10, right=1.5cm of cells] (beads) {Gel Beads\\(barcoded)};
\node[boxstep, fill=gray!10, below=0.8cm of cells] (oil) {Partitioning\\oil};

% GEM formation
\node[gem, right=1.5cm of oil] (gem) {GEM\\Cell + Bead\\+ RT reagents};

% Arrow to GEM
\draw[arrow] (cells.south east) -- (gem.north west);
\draw[arrow] (beads.south west) -- (gem.north east);
\draw[arrow] (oil.east) -- (gem.west);

% Inside GEM: cell lysis + RT
\node[boxstep, fill=green!10, right=2.0cm of gem] (rt) {Cell lysis\\+ Reverse\\transcription};
\draw[arrow] (gem.east) -- (rt.west);

% cDNA with barcodes
\node[boxstep, fill=red!10, right=1.5cm of rt] (cdna) {Barcoded\\cDNA};
\draw[arrow] (rt.east) -- (cdna.west);

% Library prep
\node[boxstep, fill=purple!10, below=1.2cm of cdna] (lib) {Library\\construction\\+ Sequencing};
\draw[arrow] (cdna.south) -- (lib.north);

% Barcode structure annotation
\node[below=1.5cm of gem, draw, rounded corners, fill=white, minimum width=9cm, minimum height=1.2cm] (barcode) {};
\node[font=\scriptsize\ttfamily, anchor=center] at (barcode.center) {
  \textcolor{BrickRed}{Illumina R1} ---
  \textcolor{blue}{Cell Barcode (16\,nt)} ---
  \textcolor{OliveGreen}{UMI (12\,nt)} ---
  \textcolor{black}{cDNA insert} ---
  \textcolor{BrickRed}{R2}
};
\node[label, above=0.05cm of barcode] {Structure of a barcoded cDNA molecule:};

\end{tikzpicture}
\caption{Schematic of the 10x Genomics Chromium droplet-based scRNA-seq
workflow. Single cells and gel beads are co-encapsulated in Gel Bead-in-Emulsions
(GEMs), where cell lysis and reverse transcription occur. Each gel bead carries
millions of oligonucleotide copies with a shared cell barcode and diverse UMIs.
The resulting cDNA is uniquely labeled with a cell barcode and a UMI for each
original mRNA molecule.}
\label{fig:scrnaseq:10x}
\end{figure}

\subsection{Gel Bead-in-Emulsions (GEMs)}

The core of the 10x Chromium system is the partitioning of single cells with
gel beads into nanoliter aqueous droplets suspended in oil (GEMs):

\begin{enumerate}[nosep]
  \item A microfluidic chip combines three streams: cell suspension, gel bead
    suspension, and partitioning oil.
  \item Each GEM ideally contains one cell and one gel bead. The vast majority
    of GEMs ($>$90\%) contain zero cells (empty) to minimize doublets (GEMs
    with two or more cells).
  \item The gel bead dissolves, releasing oligonucleotides into the droplet.
  \item Cell lysis releases mRNA, which is captured by oligo-dT sequences on
    the bead-derived oligonucleotides.
  \item Reverse transcription occurs within the GEM, producing first-strand
    cDNA tagged with the cell barcode and a unique molecular identifier (UMI).
\end{enumerate}

\subsection{Cell Barcodes and UMIs}

\begin{keyconcept}[title={Cell Barcodes and UMIs}]
\begin{description}[nosep]
  \item[Cell barcode (16\,nt):] Identifies which cell a cDNA molecule came from.
    All oligonucleotides on a single gel bead share the same cell barcode.
    Different gel beads have different barcodes. The 10x barcode whitelist
    contains $\sim$737,000 valid sequences (for 3$'$ v3/v3.1 chemistry).
  \item[UMI (12\,nt):] Unique Molecular Identifier. Each oligonucleotide on a
    bead has a random 12-mer UMI, enabling distinction between PCR duplicates
    and distinct original mRNA molecules. After PCR amplification, reads sharing
    the same cell barcode, UMI, and gene are collapsed to a single count,
    eliminating amplification bias.
\end{description}
\end{keyconcept}

\subsection{3$'$ vs.\ 5$'$ Gene Expression Chemistry}

\begin{description}[style=nextline]
  \item[3$'$ Gene Expression (GEX)] Captures the 3$'$ end of polyadenylated
    transcripts (within $\sim$500\,bp of the poly-A tail). Most commonly used
    chemistry. Sufficient for gene-level quantification and cell type
    identification. Cannot provide full-length transcript or isoform information.
  \item[5$'$ Gene Expression (GEX)] Captures the 5$'$ end of transcripts using
    template-switching. Enables simultaneous profiling of gene expression and
    immune receptor sequences (TCR/BCR V(D)J repertoires) from the same cells.
    Preferred for immunology studies.
\end{description}

\subsection{Cell Ranger Pipeline}

\software{Cell Ranger} is the standard 10x Genomics analysis pipeline that
processes raw Illumina BCL files into a gene-by-cell count matrix:

\begin{enumerate}[nosep]
  \item \texttt{cellranger mkfastq}: Demultiplexes BCL files into FASTQ files.
  \item \texttt{cellranger count}: Performs alignment (STAR), barcode counting,
    UMI deduplication, and generates the filtered feature-barcode matrix.
  \item \texttt{cellranger aggr}: Aggregates multiple \texttt{count} runs and
    normalizes across samples.
\end{enumerate}

The key output is the \textbf{filtered feature-barcode matrix}: a sparse
matrix where rows are genes ($\sim$33,000 for human) and columns are cell
barcodes that passed the cell-calling algorithm (EmptyDrops-like knee/inflection
point detection).

\subsection{Typical Throughput and Sequencing Recommendations}

\begin{itemize}[nosep]
  \item \textbf{Cells per reaction:} 500--10,000 (recommended loading for
    $\sim$0.8\% multiplet rate per 1,000 cells).
  \item \textbf{Sequencing depth:} 20,000--50,000 reads per cell for 3$'$ GEX.
    Higher depth yields diminishing returns for gene detection beyond
    $\sim$50,000 reads/cell.
  \item \textbf{Read configuration:} Read 1: 28\,bp (16\,bp cell barcode +
    12\,bp UMI); Index 1: 10\,bp (sample index); Read 2: 90\,bp (cDNA insert).
  \item \textbf{Expected genes per cell:} 2,000--6,000 median genes/cell,
    depending on cell type and sequencing depth.
\end{itemize}

\subsection{10x Chromium Flex (Fixed RNA Profiling)}

The Flex assay uses a probe-based approach to capture gene expression from
\emph{fixed} cells:

\begin{itemize}[nosep]
  \item Cells are fixed with paraformaldehyde, enabling sample banking and
    batch processing.
  \item Probe pairs hybridize to target mRNAs inside fixed cells; probes are
    ligated and then captured during GEM partitioning.
  \item Compatible with FFPE tissue sections (after deparaffinization).
  \item Supports sample multiplexing: up to 128 samples pooled in one Chromium
    run using probe-barcode-based demultiplexing.
  \item Lower background noise compared to 3$'$ GEX because probe-based
    detection avoids ambient RNA contamination.
\end{itemize}


% ---------------------------------------------------------------------------
\section{Key Computational Challenges}
\label{sec:scrnaseq:challenges}

\subsection{Sparsity and Dropout}

scRNA-seq data is extremely sparse: the median gene detection rate per cell is
only 15--30\% of expressed genes. Many genes that are truly expressed in a cell
are recorded as zero counts---a phenomenon called \textbf{dropout}. Dropout
arises from:

\begin{itemize}[nosep]
  \item Low mRNA capture efficiency ($\sim$5--15\% of transcripts in a cell for
    droplet-based methods).
  \item Stochastic sampling: with only $\sim$5,000--10,000 unique mRNA molecules
    detected per cell, lowly expressed genes are frequently missed.
  \item Shallow sequencing: not all cDNA molecules in the library are sequenced.
\end{itemize}

\begin{warningbox}
Dropout means that a zero count in scRNA-seq does not necessarily mean the gene
is not expressed. This has profound implications for analysis: methods that
treat zeros as true absence can produce misleading results. Many specialized
scRNA-seq statistical methods (e.g., MAST, scVI) explicitly model dropout or
use zero-inflated distributions.
\end{warningbox}

\subsection{Doublet Detection}

Doublets (or multiplets) occur when two or more cells are captured in the same
droplet or well. They appear as single ``cells'' with hybrid transcriptomes
that can create spurious intermediate populations or bridge clusters
artifactually.

The expected doublet rate for 10x Chromium is $\sim$0.8\% per 1,000 cells
loaded (e.g., loading 10,000 cells yields $\sim$8\% doublets).

\begin{description}[style=nextline]
  \item[\software{Scrublet}] Simulates doublets by combining random pairs of
    observed cell profiles and scores each real cell by its similarity to
    simulated doublets. Python-based.
  \item[\software{DoubletFinder}] Similar simulation-based approach implemented
    in R/Seurat. Requires parameter optimization (pK, expected doublet rate).
  \item[\software{scDblFinder}] Fast and accurate doublet detection using a
    gradient-boosted classifier trained on simulated doublets.
    R/Bioconductor-based.
\end{description}

\subsubsection{Doublet Detection: Principles and Expected Rates}
\label{subsec:scrnaseq:doublet_principles}

The fundamental principle behind simulation-based doublet detectors such as
\software{Scrublet} is straightforward:

\begin{enumerate}[nosep]
  \item \textbf{Simulate artificial doublets.} Randomly select pairs of observed
    cell profiles from the count matrix and sum their expression vectors. This
    creates synthetic doublets that mimic the transcriptomes of two co-captured
    cells.
  \item \textbf{Embed jointly.} Combine the real cells and simulated doublets,
    then perform dimensionality reduction (PCA followed by KNN graph
    construction) on the merged dataset.
  \item \textbf{Score each real cell.} For each observed cell, compute a
    \emph{doublet score} based on the local density of simulated doublets in
    its neighborhood. Cells surrounded by many simulated doublets are likely
    real doublets.
  \item \textbf{Threshold.} Apply a bimodal threshold (or a user-specified
    cutoff) to classify cells as singlets or doublets.
\end{enumerate}

The expected doublet rate for droplet-based platforms follows an approximately
linear relationship with loading density:
\begin{equation}
  \label{eq:scrnaseq:doubletrate}
  r_d \approx 0.008 \times \frac{n_{\text{cells}}}{1000}
\end{equation}
where $n_{\text{cells}}$ is the number of recovered cells per reaction. For
example, loading 5,000 cells yields an expected doublet rate of $\sim$4\%, while
10,000 cells gives $\sim$8\%. This formula assumes Poisson loading statistics
and is calibrated for the 10x Chromium platform.

\begin{warningbox}
Simulation-based doublet detectors excel at identifying \textbf{heterotypic
doublets} (two cells of different types, e.g., a T cell + a monocyte) because
these produce distinctive hybrid profiles. \textbf{Homotypic doublets} (two
cells of the same type) are nearly impossible to detect computationally because
their transcriptomes resemble a single cell with higher UMI counts. Homotypic
doublets inflate the apparent cell counts but typically do not create artifactual
clusters. When reporting cell numbers, account for both detected and undetected
(homotypic) doublets.
\end{warningbox}

\subsection{Ambient RNA Contamination}

During droplet generation, free-floating mRNA from lysed or damaged cells
(``soup'') is present in the cell suspension. This ambient RNA is captured in
every GEM, contaminating the transcriptomes of all cells with a background
signature (often enriched for highly expressed genes like hemoglobin in blood
or mitochondrial genes from damaged cells).

\begin{description}[style=nextline]
  \item[\software{SoupX}] Estimates the contamination fraction using the
    expression of genes expected to be cell-type-specific (e.g., hemoglobin in
    non-erythrocyte clusters). Subtracts estimated ambient counts.
  \item[\software{CellBender} (\texttt{remove-background})] Deep
    generative model that distinguishes cell-containing droplets from empty
    droplets and removes ambient RNA signal. Computationally intensive but
    highly effective.
  \item[\software{DecontX}] Bayesian method for ambient RNA decontamination
    that models each cell's expression as a mixture of its native profile and
    the ambient contamination profile.
\end{description}


% ---------------------------------------------------------------------------
\section{scRNA-seq Analysis Workflow}
\label{sec:scrnaseq:workflow}

The computational analysis of scRNA-seq data involves a well-established
sequence of steps, from raw data processing to biological interpretation.

% TikZ diagram: scRNA-seq analysis workflow
\begin{figure}[htbp]
\centering
\begin{tikzpicture}[
  node distance=0.65cm,
  every node/.style={font=\small},
  boxstep/.style={draw, rounded corners, fill=MidnightBlue!10, minimum width=4.5cm, minimum height=0.7cm, align=center},
  tool/.style={font=\scriptsize\ttfamily, text=OliveGreen!80!black},
  arrow/.style={-{Stealth[length=5pt]}, thick, MidnightBlue!60},
]

\node[boxstep] (raw) {Raw BCL / FASTQ};
\node[boxstep, below=of raw] (count) {Cell barcode demux + alignment + UMI counting};
\node[boxstep, below=of count] (qc) {QC filtering (nGene, nUMI, \%MT, doublets)};
\node[boxstep, below=of qc] (norm) {Normalization + HVG selection};
\node[boxstep, below=of norm] (pca) {PCA (dimensionality reduction)};
\node[boxstep, below=of pca] (cluster) {KNN graph + Leiden clustering};
\node[boxstep, below=of cluster] (umap) {UMAP / t-SNE visualization};
\node[boxstep, below=of umap] (annot) {Cell type annotation};

% Downstream branches
\node[boxstep, below left=1.0cm and 0.5cm of annot, minimum width=3.2cm] (de) {Differential\\expression};
\node[boxstep, below=1.0cm of annot, minimum width=3.2cm] (traj) {Trajectory /\\pseudotime};
\node[boxstep, below right=1.0cm and 0.5cm of annot, minimum width=3.2cm] (comm) {Cell--cell\\communication};

% Arrows
\draw[arrow] (raw) -- (count);
\draw[arrow] (count) -- (qc);
\draw[arrow] (qc) -- (norm);
\draw[arrow] (norm) -- (pca);
\draw[arrow] (pca) -- (cluster);
\draw[arrow] (cluster) -- (umap);
\draw[arrow] (umap) -- (annot);
\draw[arrow] (annot) -| (de);
\draw[arrow] (annot) -- (traj);
\draw[arrow] (annot) -| (comm);

% Tool annotations
\node[tool, right=0.15cm of count] {Cell Ranger, STARsolo, alevin-fry};
\node[tool, right=0.15cm of qc] {Scrublet, SoupX, CellBender};
\node[tool, right=0.15cm of norm] {SCTransform, scran, scVI};
\node[tool, right=0.15cm of cluster] {Scanpy, Seurat};
\node[tool, right=0.15cm of annot] {SingleR, CellTypist, Azimuth};
\node[tool, below=0.05cm of de] {\scriptsize\ttfamily DESeq2 (pseudobulk)};
\node[tool, below=0.05cm of traj] {\scriptsize\ttfamily Monocle3, scVelo};
\node[tool, below=0.05cm of comm] {\scriptsize\ttfamily CellChat, NicheNet};

\end{tikzpicture}
\caption{Standard computational workflow for scRNA-seq analysis. After
generating a cell-by-gene count matrix, cells are filtered, normalized,
dimensionality-reduced, clustered, and annotated. Downstream analyses include
differential expression, trajectory inference, and cell--cell communication.}
\label{fig:scrnaseq:workflow}
\end{figure}

\subsection{Raw Data Processing}

\begin{description}[style=nextline]
  \item[\software{Cell Ranger}] The default pipeline for 10x Chromium data (see
    \cref{sec:scrnaseq:10x}). Closed-source but free. Produces the filtered
    count matrix and a web summary report.
  \item[\software{STARsolo}] Open-source alternative to Cell Ranger, built into
    the STAR aligner. Produces equivalent results with greater flexibility and
    often faster run time. Supports 10x, Drop-seq, and other barcode schemes.
  \item[\software{alevin-fry} (Salmon)] Lightweight mapping-based quantification
    for scRNA-seq. Uses the Salmon quasi-mapping engine with barcode-aware
    processing. Produces a count matrix compatible with downstream tools. The
    \software{simpleaf} wrapper simplifies execution.
\end{description}

\subsection{Quality Control and Filtering}

After generating the raw count matrix, per-cell QC metrics are computed and
cells failing quality thresholds are removed:

\begin{itemize}[nosep]
  \item \textbf{Number of genes detected (nGene):} Cells with very few genes
    ($< 200$--500) are likely empty droplets or debris.
  \item \textbf{Number of UMIs (nUMI / nCount):} Total counts per cell.
    Extremely high values may indicate doublets.
  \item \textbf{Percentage of mitochondrial reads (\%MT):} High mitochondrial
    fraction ($> 10$--20\%, depending on tissue) indicates dying or damaged
    cells whose cytoplasmic mRNA has leaked out while mitochondrial mRNA
    (enclosed in a double membrane) is retained.
  \item \textbf{Percentage of ribosomal protein genes (\%Ribo):} Can help
    identify low-quality cells or specific cell types.
  \item \textbf{Doublet scores:} From Scrublet, DoubletFinder, or scDblFinder.
\end{itemize}

\begin{tipbox}
QC thresholds should be adapted to the tissue and experiment. Brain tissue
typically has higher mitochondrial content than PBMCs. Use distribution-based
thresholds (e.g., $> 3$ median absolute deviations from the median) rather
than fixed cutoffs when possible. The \software{scater} package in R provides
automated outlier detection with \texttt{isOutlier()}.
\end{tipbox}

\subsection{Normalization}

Raw UMI counts must be normalized to account for differences in sequencing
depth (total UMIs) across cells.

\begin{description}[style=nextline]
  \item[Log-normalization (Seurat default)] Divide each cell's counts by total
    UMIs, multiply by a scale factor (default 10,000), and log-transform:
    $\log_2(\text{counts}/\text{total} \times 10{,}000 + 1)$. Simple and
    widely used.
  \item[\software{scran} pooling-based normalization] Pools cells into groups,
    estimates pool-based size factors, and deconvolves to per-cell size factors.
    More robust than simple library-size normalization because it handles the
    sparsity of single-cell data.
  \item[\software{SCTransform} (Seurat v3+)] Fits a regularized negative
    binomial regression model to each gene as a function of total UMIs. Replaces
    both normalization and variance stabilization in a single step. Handles
    the mean-variance relationship better than log-normalization for downstream
    analyses (PCA, clustering).
  \item[\software{scVI} (scvi-tools)] Deep generative model (variational
    autoencoder) that models the count generation process, implicitly handling
    normalization, batch effects, and dropout simultaneously.
\end{description}

\subsection{Feature Selection: Highly Variable Genes}

Not all $\sim$20,000+ detected genes are informative for distinguishing cell
types. Feature selection identifies \textbf{highly variable genes (HVGs)}---those
with expression variation exceeding what is expected from sampling noise.

\begin{itemize}[nosep]
  \item Seurat: fits a variance-stabilizing transformation and selects the top
    2,000--5,000 genes by residual variance.
  \item Scanpy: uses the method of Zheng et al.\ or the Seurat-style approach
    to identify HVGs.
  \item Only HVGs are used for dimensionality reduction and clustering,
    reducing noise and computation.
\end{itemize}

\subsection{Dimensionality Reduction}

\begin{description}[style=nextline]
  \item[PCA (Principal Component Analysis)] Linear dimensionality reduction.
    Typically the first 20--50 principal components are retained, capturing
    the major axes of variation. PCA is performed on the HVG expression matrix.
  \item[t-SNE (t-distributed Stochastic Neighbor Embedding)] Non-linear
    dimensionality reduction that preserves local structure. Produces 2D
    embeddings widely used for visualization. Sensitive to hyperparameters
    (perplexity). Does not preserve global distances.
  \item[UMAP (Uniform Manifold Approximation and Projection)] Non-linear
    method that preserves both local and (approximately) global structure.
    Faster than t-SNE and currently the most popular visualization method
    for scRNA-seq. Operates on the PCA-reduced space (not raw expression).
\end{description}

\begin{warningbox}
t-SNE and UMAP are visualization tools, not analytical methods. Do not perform
clustering on t-SNE/UMAP coordinates. Cluster on the PCA space or the shared
nearest neighbor (SNN) graph. Distances and apparent cluster sizes in t-SNE/UMAP
plots can be misleading.
\end{warningbox}

\subsection{Mathematical Foundations of t-SNE and UMAP}
\label{subsec:scrnaseq:tsne_umap_math}

Although t-SNE and UMAP are often treated as black-box visualization tools,
understanding their mathematical principles is essential for interpreting their
output correctly and tuning hyperparameters.

\subsubsection{t-SNE: t-distributed Stochastic Neighbor Embedding}

t-SNE operates in two stages. First, in the high-dimensional PCA space, pairwise
similarities are converted into conditional probabilities. The similarity of
point $x_j$ to point $x_i$ is defined as:
\begin{equation}
  \label{eq:scrnaseq:tsne_pij}
  p_{j|i} = \frac{\exp\!\bigl(-\|x_i - x_j\|^2 / 2\sigma_i^2\bigr)}
    {\sum_{k \neq i} \exp\!\bigl(-\|x_i - x_k\|^2 / 2\sigma_i^2\bigr)}
\end{equation}
where $\sigma_i$ is chosen such that the perplexity of the conditional
distribution matches a user-specified value (typically 30). The perplexity
controls the effective number of neighbors considered for each point.

The joint probability matrix is symmetrized as
$p_{ij} = (p_{j|i} + p_{i|j}) / (2n)$. In the low-dimensional embedding
(2D), similarities are modeled using a Student's $t$-distribution with one
degree of freedom (Cauchy distribution):
\begin{equation}
  \label{eq:scrnaseq:tsne_qij}
  q_{ij} = \frac{(1 + \|y_i - y_j\|^2)^{-1}}
    {\sum_{k \neq l} (1 + \|y_k - y_l\|^2)^{-1}}
\end{equation}
where $y_i, y_j$ are the low-dimensional coordinates. The heavy tails of the
$t$-distribution allow nearby points to be modeled faithfully while pushing
dissimilar points further apart, alleviating the ``crowding problem'' that
affects Gaussian-based methods like classical SNE.

t-SNE minimizes the Kullback--Leibler divergence between $P$ and $Q$:
\begin{equation}
  \label{eq:scrnaseq:tsne_kl}
  KL(P \| Q) = \sum_i \sum_{j \neq i} p_{ij}
    \log \frac{p_{ij}}{q_{ij}}
\end{equation}
via gradient descent. The asymmetry of KL divergence means that t-SNE penalizes
placing similar points far apart (``tearing'' neighborhoods) more heavily than
placing dissimilar points close together. This is why t-SNE excels at preserving
local structure but can distort global distances.

\subsubsection{UMAP: Uniform Manifold Approximation and Projection}

UMAP is grounded in topological data analysis and Riemannian geometry. Its key
steps are:

\begin{enumerate}[nosep]
  \item \textbf{Fuzzy simplicial set construction.} For each point $x_i$, UMAP
    finds its $k$ nearest neighbors and converts distances into membership
    strengths using a smooth kernel:
    \begin{equation}
      \label{eq:scrnaseq:umap_weight}
      w_{ij} = \exp\!\left(-\frac{d(x_i, x_j) - \rho_i}{\sigma_i}\right)
    \end{equation}
    where $\rho_i$ is the distance to the nearest neighbor of $x_i$ (ensuring
    local connectivity) and $\sigma_i$ is scaled to match the desired number of
    neighbors. These weights are symmetrized to form a fuzzy simplicial set
    (weighted graph) representing the topological structure of the data.
  \item \textbf{Low-dimensional optimization.} UMAP initializes a low-dimensional
    embedding (typically using spectral initialization on the fuzzy graph
    Laplacian) and optimizes it by minimizing a \textbf{cross-entropy loss}
    between the high-dimensional and low-dimensional fuzzy sets:
    \begin{equation}
      \label{eq:scrnaseq:umap_loss}
      \mathcal{L} = \sum_{(i,j)} \Bigl[
        w_{ij} \log\!\frac{w_{ij}}{v_{ij}}
        + (1 - w_{ij}) \log\!\frac{1 - w_{ij}}{1 - v_{ij}}
      \Bigr]
    \end{equation}
    where $v_{ij} = (1 + a \|y_i - y_j\|^{2b})^{-1}$ is the low-dimensional
    similarity, with $a$ and $b$ determined by the \texttt{min\_dist} parameter.
  \item \textbf{Stochastic gradient descent with negative sampling.} Rather than
    computing all pairwise interactions, UMAP uses negative sampling to
    approximate the repulsive term, making it substantially faster than t-SNE
    for large datasets ($\mathcal{O}(n)$ per epoch vs.\ $\mathcal{O}(n^2)$ for
    exact t-SNE).
\end{enumerate}

\begin{keyconcept}[title={t-SNE vs.\ UMAP: Key Differences}]
\begin{itemize}[nosep]
  \item \textbf{Global structure:} UMAP's cross-entropy loss includes a
    repulsive term for dissimilar points, better preserving global relationships
    (relative cluster positions) compared to t-SNE's KL divergence.
  \item \textbf{Speed:} UMAP is typically 5--50$\times$ faster than t-SNE for
    datasets with $>$50,000 cells, due to negative sampling and graph-based
    initialization.
  \item \textbf{Reproducibility:} UMAP with fixed random seed produces
    consistent embeddings; t-SNE results can vary substantially between runs.
  \item \textbf{Key hyperparameters:} For t-SNE, perplexity (30--50 typical);
    for UMAP, \texttt{n\_neighbors} (15--30) and \texttt{min\_dist}
    (0.1--0.5). Both should be explored rather than left at defaults.
\end{itemize}
\end{keyconcept}

\subsection{Clustering}

Graph-based clustering is the standard approach for scRNA-seq:

\begin{enumerate}[nosep]
  \item Construct a k-nearest neighbor (KNN) graph in PCA space. Each cell is a
    node; edges connect cells to their $k$ nearest neighbors.
  \item Refine edge weights using shared nearest neighbors (SNN graph) or
    Jaccard similarity.
  \item Apply a community detection algorithm:
    \begin{itemize}[nosep]
      \item \textbf{Louvain algorithm:} Optimizes modularity; fast but can
        produce poorly connected clusters.
      \item \textbf{Leiden algorithm:} Improved version of Louvain that
        guarantees well-connected communities. Generally preferred.
    \end{itemize}
  \item The \texttt{resolution} parameter controls the granularity of
    clustering: higher values produce more clusters. There is no single
    ``correct'' resolution; it should be guided by biological knowledge.
\end{enumerate}

\subsection{Cell Type Annotation}

Assigning biological identity to clusters is one of the most critical---and
most subjective---steps in scRNA-seq analysis.

\begin{description}[style=nextline]
  \item[Manual annotation] Examine marker genes for each cluster (e.g.,
    \gene{CD3D}/\gene{CD3E} for T cells, \gene{MS4A1}/\gene{CD79A} for B
    cells, \gene{LYZ}/\gene{CD14} for monocytes). Requires domain expertise.
    The \texttt{FindAllMarkers} function in Seurat and
    \texttt{sc.tl.rank\_genes\_groups} in Scanpy identify cluster-specific
    marker genes.
  \item[\software{SingleR}] Automated annotation by correlating each cell's
    expression profile to reference transcriptomic datasets (e.g., Human Primary
    Cell Atlas, Monaco Immune Data, Blueprint Encode). R-based.
  \item[\software{CellTypist}] Machine-learning-based classifier trained on
    large reference datasets. Pre-trained models for immune cells, brain cells,
    and pan-tissue atlases. Python-based.
  \item[\software{scType}] Fully automated cell type annotation using
    curated marker gene databases. R-based; no reference dataset required.
  \item[\software{Azimuth}] Web-based and R-based tool from the Seurat
    developers. Uses reference mapping to project query cells onto well-curated
    reference atlases (PBMC, lung, motor cortex, etc.). Provides both coarse
    and fine-grained cell type labels with confidence scores.
\end{description}

\subsection{Differential Expression}

Identifying genes that are differentially expressed between clusters or
conditions requires methods appropriate for the distributional properties of
scRNA-seq data:

\begin{description}[style=nextline]
  \item[Wilcoxon rank-sum test] Non-parametric; the default in both Seurat and
    Scanpy. Fast and performs reasonably well for cluster markers but does not
    account for the hierarchical structure of cells within samples.
  \item[\software{MAST}] Hurdle model that jointly models the rate of
    expression (logistic regression for zero vs.\ non-zero) and the expression
    level (linear model for non-zero values). Accounts for cellular detection
    rate as a confounding variable.
  \item[Pseudobulk approaches] Aggregate counts per cell type per biological
    replicate, then apply bulk DE tools (DESeq2, edgeR). This approach correctly
    treats biological replicates as the unit of inference (not individual
    cells) and has been shown to have the best calibrated false discovery rates
    in recent benchmarks. \textbf{Recommended for comparing conditions (e.g.,
    treatment vs.\ control) when biological replicates are available.}
\end{description}

\begin{keyconcept}[title={The Pseudobulk Principle}]
In a scRNA-seq experiment with multiple biological samples, individual cells
are not independent observations---cells from the same sample share the same
biological variation. Treating each cell as an independent replicate inflates
the sample size artificially and leads to extreme $p$-value inflation. The
solution is pseudobulk aggregation: sum counts across cells of the same type
within each biological replicate, then use bulk DE methods on the resulting
replicate-level counts.
\end{keyconcept}

\subsection{Trajectory Analysis and Pseudotime}

Many biological processes (differentiation, activation, disease progression)
involve continuous changes in gene expression rather than discrete states.
Trajectory analysis orders cells along a ``pseudotime'' axis that recapitulates
the dynamic process.

\begin{description}[style=nextline]
  \item[\software{Monocle3}] Learns a principal graph (trajectory) through the
    UMAP embedding using reversed graph embedding. Assigns each cell a
    pseudotime value representing its position along the trajectory. Supports
    branching trajectories for divergent differentiation paths.
  \item[\software{Slingshot}] Fits smooth curves through PCA or diffusion map
    space. Uses cluster information to define trajectory start and end points.
    Well-suited for simpler, tree-like trajectories.
  \item[\software{scVelo}] Infers RNA velocity---the time derivative of gene
    expression---from the ratio of spliced to unspliced mRNA counts in each
    cell. RNA velocity vectors indicate the future transcriptional state of
    each cell, providing directionality to trajectories.
  \item[\software{CellRank}] Combines RNA velocity with expression similarity
    to compute transition probabilities between cells. Identifies initial and
    terminal states and the most likely differentiation fates using Markov chain
    analysis.
\end{description}

\subsubsection{RNA Velocity: Mathematical Framework}
\label{subsec:scrnaseq:rna_velocity}

RNA velocity exploits the fact that standard scRNA-seq protocols capture both
spliced (mature) and unspliced (intron-containing) mRNA molecules. The ratio
between unspliced and spliced counts provides information about the
\emph{direction} of transcriptional change for each gene in each cell.

The core dynamical model for a single gene describes the time evolution of
unspliced ($u$) and spliced ($s$) mRNA abundances:
\begin{align}
  \label{eq:scrnaseq:velocity_unspliced}
  \frac{du}{dt} &= \alpha(t) - \beta\, u \\
  \label{eq:scrnaseq:velocity_spliced}
  \frac{ds}{dt} &= \beta\, u - \gamma\, s
\end{align}
where $\alpha(t)$ is the transcription rate, $\beta$ is the splicing rate
(conversion of unspliced to spliced mRNA), and $\gamma$ is the degradation rate
of spliced mRNA. The RNA \emph{velocity} for each gene is defined as $ds/dt$,
representing the expected rate of change of the mature mRNA.

\begin{description}[style=nextline]
  \item[Steady-state model (\software{velocyto})] Assumes that genes are in
    either an induction or repression phase, each at approximate steady state.
    At steady state, $ds/dt = 0$ implies $\beta u = \gamma s$, so the expected
    ratio is $u/s = \gamma/\beta$. Cells deviating above this ratio (excess
    unspliced) are actively upregulating the gene; cells below are downregulating
    it. The velocity is estimated as:
    \begin{equation}
      \label{eq:scrnaseq:velocity_ss}
      v_g = \beta\, u_g - \gamma_g\, s_g
    \end{equation}
    where $\gamma_g$ is estimated from the slope of the upper quantile of the
    $u$ vs.\ $s$ phase portrait.
  \item[Dynamical model (\software{scVelo})] Relaxes the steady-state assumption
    by solving the full system of ODEs (Equations
    \ref{eq:scrnaseq:velocity_unspliced}--\ref{eq:scrnaseq:velocity_spliced})
    using an expectation-maximization algorithm. This jointly infers the kinetic
    parameters $\alpha$, $\beta$, $\gamma$, and a latent time for each cell,
    providing a more accurate velocity field for genes that are not at steady
    state (e.g., during rapid transitions).
  \item[Multi-omic velocity (\software{MultiVelo})] Extends the velocity
    framework to incorporate chromatin accessibility (from multi-omic
    ATAC + RNA data), modeling the cascade from chromatin opening to unspliced
    to spliced mRNA.
\end{description}

\begin{tipbox}
RNA velocity estimates are gene-specific vectors in expression space. To
visualize them on a UMAP or t-SNE embedding, the high-dimensional velocity
vectors are projected onto the 2D embedding by computing the correlation between
each cell's velocity and the expression difference to its neighbors. Arrows on
the embedding indicate the predicted direction of transcriptional change, not
physical movement.
\end{tipbox}

\subsection{Cell--Cell Communication}

\begin{description}[style=nextline]
  \item[\software{CellChat}] Infers intercellular communication networks from
    scRNA-seq data using a curated database of ligand--receptor interactions
    (including multimeric complexes). Models communication probability based on
    the law of mass action and identifies significant signaling pathways between
    cell types.
  \item[\software{NicheNet}] Predicts which ligands from sender cells best
    explain gene expression changes in receiver cells, integrating prior
    knowledge of signaling and gene regulatory networks.
  \item[\software{LIANA}] A framework that benchmarks and combines multiple
    cell--cell communication tools (CellPhoneDB, CellChat, NATMI,
    SingleCellSignalR, etc.) to produce consensus interaction scores.
\end{description}

\subsection{Gene Regulatory Network Inference}

\begin{description}[style=nextline]
  \item[\software{SCENIC / pySCENIC}] \textbf{S}ingle-\textbf{C}ell r\textbf{E}gulatory
    \textbf{N}etwork \textbf{I}nference and \textbf{C}lustering. A three-step
    pipeline:
    \begin{enumerate}[nosep]
      \item Co-expression modules: identify genes co-expressed with each
        transcription factor (TF) using GRNBoost2 or GENIE3.
      \item Regulon pruning: retain only target genes with TF binding motifs
        in their regulatory regions (cis-regulatory analysis using RcisTarget
        or pycisTarget).
      \item Regulon activity scoring: quantify the activity of each regulon
        (TF + targets) in each cell using AUCell.
    \end{enumerate}
    SCENIC regulon activity provides a biologically interpretable dimensionality
    reduction and can reveal TF-driven cell states invisible in gene expression
    alone.
\end{description}


% ---------------------------------------------------------------------------
\section{Analysis Frameworks: Scanpy vs.\ Seurat}
\label{sec:scrnaseq:frameworks}

Two comprehensive frameworks dominate the scRNA-seq analysis ecosystem:

\begin{table}[htbp]
\centering
\caption{Comparison of the Scanpy (Python) and Seurat (R) scRNA-seq analysis
frameworks.}
\label{tab:scrnaseq:frameworks}
\small
\begin{tabularx}{\textwidth}{lXX}
\toprule
\textbf{Feature} & \textbf{Scanpy (Python)} & \textbf{Seurat (R)} \\
\midrule
Language & Python & R \\
Current version & Scanpy 1.10+ & Seurat v5 \\
Data structure & AnnData (\texttt{.h5ad}) & SeuratObject / BPCells \\
Core functions & Preprocessing, clustering, DE, visualization & Preprocessing, clustering, DE, visualization \\
Normalization & Log-norm, scran (via rpy2) & Log-norm, SCTransform, sketch-based \\
Integration & scVI, scANVI, BBKNN, Harmony & CCA, RPCA, Harmony, scVI bridge \\
Scalability & Excellent (sparse arrays, GPU via rapids-singlecell) & Improved in v5 (BPCells on-disk), sketch-based \\
Visualization & Matplotlib-based (highly customizable) & ggplot2-based (publication-ready) \\
Ecosystem & scvi-tools, squidpy, CellRank, decoupler & Signac, Azimuth, BPCells \\
Community & Strong in computational biology / ML & Strong in biomedical research \\
\bottomrule
\end{tabularx}
\end{table}

\begin{tipbox}
Both frameworks are excellent and actively maintained. Choose based on your
language proficiency and collaborators' preferences. Many analyses use both:
Scanpy for preprocessing and Seurat for visualization (or vice versa).
Conversion between AnnData (\texttt{.h5ad}) and SeuratObject formats is
straightforward using the \software{SeuratDisk} or \software{MuDataSeurat}
packages.
\end{tipbox}


% ---------------------------------------------------------------------------
\section{Integration of Multiple Datasets}
\label{sec:scrnaseq:integration}

Combining scRNA-seq datasets from different experiments, batches, or
technologies is essential for large-scale analyses but is complicated by
batch effects. Integration methods aim to place cells from different batches
into a common embedding while preserving biological variation.

\begin{description}[style=nextline]
  \item[\software{Harmony}] Fast linear integration that iteratively adjusts
    PCA embeddings to remove batch effects while preserving biological clusters.
    Extremely popular due to speed and simplicity. Works as a post-PCA
    correction.
  \item[\software{scVI} (Variational Inference)] Deep generative model
    (variational autoencoder) that learns a batch-corrected latent space.
    Handles complex batch effects and provides uncertainty estimates. The gold
    standard for large, complex integration tasks.
  \item[\software{BBKNN} (Batch Balanced KNN)] Modifies the KNN graph
    construction to balance neighbors across batches. Simple and effective;
    does not modify the expression matrix.
  \item[Seurat integration (CCA/RPCA)] Identifies ``anchors'' (mutual nearest
    neighbors) between datasets using canonical correlation analysis (CCA) or
    reciprocal PCA (RPCA). Anchors guide the alignment of datasets into a
    common space. RPCA is faster and recommended for datasets with partially
    overlapping cell types.
\end{description}

\begin{warningbox}
Over-integration can remove genuine biological differences between datasets.
Always verify that integration preserved known biological signals (e.g., marker
gene expression for expected cell types). Compare integrated and non-integrated
results critically, and use integration metrics (kBET, LISI, silhouette width)
to quantify batch mixing and biological conservation.
\end{warningbox}

\subsection{Batch Correction Methods: A Detailed Comparison}
\label{subsec:scrnaseq:batch_correction}

Choosing the right integration method depends on dataset size, batch complexity,
and the importance of preserving subtle biological variation. The following table
summarizes the key properties of the most widely used methods:

\begin{table}[htbp]
\centering
\caption{Comparison of batch correction and integration methods for scRNA-seq.}
\label{tab:scrnaseq:batch_correction}
\small
\begin{tabularx}{\textwidth}{lXXXX}
\toprule
\textbf{Method} & \textbf{Approach} & \textbf{Speed} & \textbf{Preservation of biology} & \textbf{Notes} \\
\midrule
\software{Harmony} & Iterative soft clustering on PCA embeddings; adjusts cluster centroids across batches & Very fast ($<$1 min for 500k cells) & Good; may over-correct rare populations & Post-PCA; does not modify expression matrix; easy to use \\
MNN (Mutual Nearest Neighbors) & Identifies mutual nearest neighbor pairs across batches; computes correction vectors & Moderate & Good for conserved populations & Implemented in \software{batchelor} (R); foundational method \\
\software{scVI} & Variational autoencoder with batch as a covariate in the generative model & Slow (GPU recommended; minutes--hours) & Excellent; learns complex non-linear corrections & Produces a latent space, not corrected counts; best for complex multi-batch designs \\
\software{BBKNN} & Constructs a batch-balanced KNN graph by trimming neighbors per batch & Very fast & Good; minimal distortion of rare types & Corrects the graph, not the embedding or expression; use with Scanpy \\
\software{Scanorama} & Panoramic stitching inspired by image alignment; finds mutual nearest neighbors in SVD space & Fast (scales to millions of cells) & Good; robust to partially overlapping datasets & Can correct expression matrix directly; useful when datasets share only some cell types \\
\bottomrule
\end{tabularx}
\end{table}

\begin{keyconcept}[title={Metrics for Evaluating Integration Quality}]
Integration quality should be assessed along two orthogonal axes:
\begin{itemize}[nosep]
  \item \textbf{Batch mixing} (technical correction): Do cells of the same type
    from different batches intermingle? Measured by kBET (k-nearest neighbor
    batch effect test) and the iLISI (integration Local Inverse Simpson Index).
    Higher values indicate better batch mixing.
  \item \textbf{Biological conservation}: Are genuine cell types still
    distinguishable after integration? Measured by cLISI (cell-type LISI),
    silhouette width, and ARI (Adjusted Rand Index) against known labels.
    Over-integration reduces these metrics.
\end{itemize}
The \software{scIB} (single-cell Integration Benchmarking) framework provides a
standardized pipeline for computing these metrics and ranking integration methods
for a given dataset.
\end{keyconcept}


% ---------------------------------------------------------------------------
\section{Large-Scale Cell Atlases}
\label{sec:scrnaseq:atlases}

The field has moved from profiling individual tissues to constructing
organism-wide cell atlases:

\begin{description}[style=nextline]
  \item[Human Cell Atlas (HCA)] An international consortium aiming to create a
    comprehensive reference map of every human cell type. Spans dozens of organs
    and developmental stages. Data accessible via the HCA Data Portal.
  \item[Tabula Sapiens] A multi-organ, single-cell transcriptomic atlas of
    \species{Homo sapiens}. Profiled $\sim$500,000 cells from 24 organs of
    individual donors, enabling cross-tissue cell type comparisons.
  \item[Tabula Muris / Tabula Muris Senis] Mouse cell atlases covering 20+
    organs using both 10x Chromium (3$'$) and Smart-seq2 (full-length) protocols.
    Senis extends to aging, covering mice from 1 to 30 months.
  \item[CellxGene Census] Aggregated and harmonized single-cell data from
    multiple studies, providing programmatic access (Python/R APIs) to tens
    of millions of cells with standardized metadata and cell type annotations.
\end{description}


% ---------------------------------------------------------------------------
\section{Practical Considerations}
\label{sec:scrnaseq:practical}

\subsection{How Many Cells to Sequence?}

\begin{itemize}[nosep]
  \item \textbf{For common cell types:} 3,000--5,000 cells per sample is often
    sufficient.
  \item \textbf{For rare populations:} If a target cell type is $\sim$1\% of
    the sample, profiling 10,000--20,000 cells ensures $\sim$100--200 cells of
    that type.
  \item \textbf{For differential abundance testing:} Multiple biological
    replicates (n $\geq$ 3) are far more important than the number of cells per
    sample.
  \item \textbf{Trade-off:} More cells at lower depth vs.\ fewer cells at
    higher depth. For cell type identification, breadth (more cells) generally
    wins; for gene-level analyses within a cell type, depth matters more.
\end{itemize}

\subsection{Sequencing Depth Per Cell}

\begin{itemize}[nosep]
  \item \textbf{3$'$ 10x Chromium:} 20,000--50,000 reads per cell is standard.
    Beyond 50,000 reads, the number of newly detected genes saturates rapidly.
  \item \textbf{5$'$ 10x Chromium:} Similar depth for GEX; additional depth for
    V(D)J libraries ($\sim$5,000 reads/cell).
  \item \textbf{Smart-seq2:} 0.5--1 million reads per cell (full-length,
    no UMIs; deeper sequencing improves isoform detection).
  \item \textbf{Multiplexed experiments:} Calculate total reads as
    (cells $\times$ reads/cell) and ensure the sequencing run provides
    sufficient output.
\end{itemize}

\subsection{Tissue Dissociation Considerations}

Preparing a high-quality single-cell suspension is often the most challenging
step in a scRNA-seq experiment:

\begin{itemize}[nosep]
  \item \textbf{Enzymatic digestion:} Collagenase, trypsin, Liberase, or
    tissue-specific cocktails. Over-digestion kills cells; under-digestion
    yields clumps and debris.
  \item \textbf{Mechanical dissociation:} Gentle trituration or the
    Miltenyi gentleMACS for standardized processing.
  \item \textbf{Dead cell removal:} MACS dead cell removal kit or FACS-based
    viability gating (DAPI, 7-AAD, calcein AM).
  \item \textbf{Dissociation-induced artifacts:} Enzymatic digestion can
    activate stress-response genes (e.g., \gene{FOS}, \gene{JUN},
    \gene{HSPA1A}). Cold-active protease protocols or fixation-based approaches
    (10x Flex, FANS for nuclei) mitigate this.
  \item \textbf{Single-nucleus RNA-seq (snRNA-seq):} For tissues that cannot be
    dissociated into viable single cells (e.g., brain, adipose, archived frozen
    tissue), nuclei are isolated and profiled. Uses the same 10x Chromium
    platform with nuclei-specific protocols. Detects both cytoplasmic and
    nuclear RNA (more intronic reads).
\end{itemize}

\begin{warningbox}
Always check for dissociation-induced gene expression artifacts by examining
immediate-early gene expression across clusters. If stress genes (e.g.,
\gene{FOS}, \gene{JUN}, \gene{ATF3}) dominate the first principal components
or separate clusters, consider regressing them out or switching to a
cold-dissociation or nucleus-based protocol.
\end{warningbox}


% ---------------------------------------------------------------------------
\section{Comparison of scRNA-seq Platforms}
\label{sec:scrnaseq:comparison}

\begin{table}[htbp]
\centering
\caption{Comparison of major scRNA-seq platforms and methods.}
\label{tab:scrnaseq:platforms}
\small
\begin{tabularx}{\textwidth}{lXXXX}
\toprule
\textbf{Feature} & \textbf{10x Chromium 3$'$} & \textbf{Smart-seq2/3} & \textbf{Drop-seq} & \textbf{sci-RNA-seq3} \\
\midrule
Isolation & Droplet & FACS to plate & Droplet & Combinatorial indexing \\
Throughput & 500--10,000 cells & 96--384 cells & 2,000--10,000 cells & $>$100,000 cells \\
Sensitivity & 3,000--6,000 genes/cell & 5,000--10,000 genes/cell & 1,500--3,000 genes/cell & 1,000--2,000 genes/cell \\
Coverage & 3$'$ end & Full-length & 3$'$ end & 3$'$ end \\
UMIs & Yes (12\,nt) & No (SS2) / Yes (SS3) & Yes & Yes \\
Cost per cell & \$0.10--0.50 & \$5--20 & \$0.05--0.20 & \$0.01--0.05 \\
Equipment & Chromium Controller & FACS + thermocycler & Custom microfluidics & Standard lab equipment \\
Isoform info & No & Yes & No & No \\
Multiplexing & CellPlex, Flex & Limited & Limited & Built-in \\
Best use case & Standard profiling & Deep per-cell profiling & Budget-friendly scaling & Organism-scale atlases \\
\bottomrule
\end{tabularx}
\end{table}


% ---------------------------------------------------------------------------
\section{Snakemake Workflow for scRNA-seq}
\label{sec:scrnaseq:snakemake}

\begin{lstlisting}[language=Python, caption={Snakemake workflow for 10x
Chromium scRNA-seq analysis using STARsolo and Scanpy.},
label={lst:scrnaseq:snakemake}, float=htbp]
# Snakefile for scRNA-seq Analysis (10x Chromium)
# ================================================
# Pipeline: STARsolo -> Scanpy (QC, norm, cluster, annotate)

import pandas as pd

configfile: "config/config.yaml"

samples = pd.read_csv(config["samples"], sep="\t")
SAMPLES = samples["sample"].tolist()

STAR_INDEX = config["star_index"]
WHITELIST  = config["barcode_whitelist"]

rule all:
    input:
        expand("results/starsolo/{sample}/Solo.out/"
               "Gene/filtered/matrix.mtx",
               sample=SAMPLES),
        "results/scanpy/integrated_analysis.h5ad",
        "results/scanpy/umap_celltypes.pdf",
        "results/multiqc/multiqc_report.html"

rule starsolo:
    input:
        r1="data/raw/{sample}_R1.fastq.gz",
        r2="data/raw/{sample}_R2.fastq.gz",
        index=STAR_INDEX,
        whitelist=WHITELIST
    output:
        bam="results/starsolo/{sample}/"
            "Aligned.sortedByCoord.out.bam",
        matrix="results/starsolo/{sample}/Solo.out/"
               "Gene/filtered/matrix.mtx",
        barcodes="results/starsolo/{sample}/Solo.out/"
                 "Gene/filtered/barcodes.tsv",
        features="results/starsolo/{sample}/Solo.out/"
                 "Gene/filtered/features.tsv",
        log="results/starsolo/{sample}/Log.final.out"
    threads: 12
    resources:
        mem_mb=40000
    params:
        prefix="results/starsolo/{sample}/",
        cb_len=16,
        umi_len=12
    conda: "envs/starsolo.yaml"
    shell:
        """
        STAR --runMode alignReads \
          --genomeDir {input.index} \
          --readFilesIn {input.r2} {input.r1} \
          --readFilesCommand zcat \
          --soloType CB_UMI_Simple \
          --soloCBwhitelist {input.whitelist} \
          --soloCBlen {params.cb_len} \
          --soloUMIlen {params.umi_len} \
          --soloFeatures Gene \
          --soloMultiMappers EM \
          --soloCellFilter EmptyDrops_CR \
          --outSAMtype BAM SortedByCoordinate \
          --outSAMattributes NH HI nM AS CR UR CB UB \
            GX GN \
          --outFileNamePrefix {params.prefix} \
          --runThreadN {threads} \
          --limitBAMsortRAM 30000000000
        """

rule doublet_detection:
    input:
        matrix="results/starsolo/{sample}/Solo.out/"
               "Gene/filtered/matrix.mtx"
    output:
        scores="results/doublets/{sample}_scores.csv"
    conda: "envs/scanpy.yaml"
    script:
        "scripts/run_scrublet.py"

rule scanpy_analysis:
    input:
        matrices=expand(
            "results/starsolo/{sample}/Solo.out/"
            "Gene/filtered/matrix.mtx",
            sample=SAMPLES),
        doublets=expand(
            "results/doublets/{sample}_scores.csv",
            sample=SAMPLES),
        sample_sheet=config["samples"]
    output:
        h5ad="results/scanpy/integrated_analysis.h5ad",
        umap="results/scanpy/umap_celltypes.pdf",
        markers="results/scanpy/marker_genes.csv",
        qc_report="results/scanpy/qc_violin_plots.pdf"
    conda: "envs/scanpy.yaml"
    script:
        "scripts/run_scanpy.py"

rule multiqc:
    input:
        expand("results/starsolo/{sample}/Log.final.out",
               sample=SAMPLES)
    output:
        "results/multiqc/multiqc_report.html"
    conda: "envs/qc.yaml"
    shell:
        "multiqc results/ -o results/multiqc/ --force"
\end{lstlisting}


% ---------------------------------------------------------------------------
\section{Nextflow Workflow for scRNA-seq}
\label{sec:scrnaseq:nextflow}

\begin{lstlisting}[language=Java, caption={Nextflow DSL2 workflow for 10x
Chromium scRNA-seq analysis using STARsolo and Scanpy.},
label={lst:scrnaseq:nextflow}, float=htbp]
#!/usr/bin/env nextflow
// scRNA-seq Analysis Pipeline (Nextflow DSL2)
// =============================================
// Pipeline: STARsolo -> Scrublet -> Scanpy

nextflow.enable.dsl = 2

params.reads        = "data/raw/*_{R1,R2}.fastq.gz"
params.star_index   = "resources/star_index"
params.whitelist    = "resources/3M-february-2018.txt"
params.sample_sheet = "config/samples.tsv"
params.outdir       = "results"

log.info """
=================================
 scRNA-seq Analysis Pipeline
=================================
 reads        : ${params.reads}
 STAR index   : ${params.star_index}
 whitelist    : ${params.whitelist}
 outdir       : ${params.outdir}
=================================
"""

process STARSOLO {
    tag "$sample_id"
    publishDir "${params.outdir}/starsolo/${sample_id}",
        mode: 'copy'
    cpus 12
    memory '40 GB'

    input:
    tuple val(sample_id), path(reads)
    path(star_index)
    path(whitelist)

    output:
    tuple val(sample_id),
        path("Solo.out/Gene/filtered/"),
        emit: filtered
    path("Log.final.out"), emit: log
    tuple val(sample_id),
        path("Aligned.sortedByCoord.out.bam"),
        emit: bam

    script:
    def (r1, r2) = reads
    """
    STAR --runMode alignReads \
      --genomeDir ${star_index} \
      --readFilesIn ${r2} ${r1} \
      --readFilesCommand zcat \
      --soloType CB_UMI_Simple \
      --soloCBwhitelist ${whitelist} \
      --soloCBlen 16 \
      --soloUMIlen 12 \
      --soloFeatures Gene \
      --soloMultiMappers EM \
      --soloCellFilter EmptyDrops_CR \
      --outSAMtype BAM SortedByCoordinate \
      --outSAMattributes NH HI nM AS CR UR CB UB \
        GX GN \
      --runThreadN ${task.cpus}
    """
}

process SCRUBLET {
    tag "$sample_id"
    publishDir "${params.outdir}/doublets", mode: 'copy'

    input:
    tuple val(sample_id), path(filtered_dir)

    output:
    tuple val(sample_id),
        path("${sample_id}_doublet_scores.csv"),
        emit: scores

    script:
    """
    python ${projectDir}/bin/run_scrublet.py \
      --input ${filtered_dir} \
      --output ${sample_id}_doublet_scores.csv
    """
}

process SCANPY_ANALYSIS {
    publishDir "${params.outdir}/scanpy", mode: 'copy'
    cpus 4
    memory '16 GB'

    input:
    path(filtered_dirs)
    path(doublet_scores)
    path(sample_sheet)

    output:
    path("integrated_analysis.h5ad"), emit: h5ad
    path("umap_celltypes.pdf"),       emit: umap
    path("marker_genes.csv"),         emit: markers
    path("qc_violin_plots.pdf"),      emit: qc

    script:
    """
    python ${projectDir}/bin/run_scanpy.py \
      --input_dirs ${filtered_dirs} \
      --doublets ${doublet_scores} \
      --samples ${sample_sheet} \
      --outdir .
    """
}

process MULTIQC {
    publishDir "${params.outdir}/multiqc", mode: 'copy'

    input:
    path('*')

    output:
    path("multiqc_report.html")

    script:
    """
    multiqc . -o . --force
    """
}

workflow {
    reads_ch = Channel
        .fromFilePairs(params.reads, checkIfExists: true)

    star_index = Channel.fromPath(params.star_index)
    whitelist  = Channel.fromPath(params.whitelist)

    // Align and quantify
    STARSOLO(reads_ch, star_index.collect(),
             whitelist.collect())

    // Detect doublets per sample
    SCRUBLET(STARSOLO.out.filtered)

    // Integrated analysis
    SCANPY_ANALYSIS(
        STARSOLO.out.filtered
            .map { it[1] }.collect(),
        SCRUBLET.out.scores
            .map { it[1] }.collect(),
        Channel.fromPath(params.sample_sheet)
    )

    // QC report
    MULTIQC(STARSOLO.out.log.collect())
}
\end{lstlisting}

\begin{tipbox}
For production scRNA-seq analysis, the \software{nf-core/scrnaseq} pipeline
provides a comprehensive Nextflow workflow supporting Cell Ranger, STARsolo,
alevin-fry (Salmon), and kallisto-bustools as quantification backends, with
automatic whitelist selection and extensive QC reporting.
\end{tipbox}


% ---------------------------------------------------------------------------
\section{Chapter Summary}
\label{sec:scrnaseq:summary}

Single-cell RNA sequencing has revolutionized our ability to dissect cellular
heterogeneity and has become an indispensable tool in modern biology. In this
chapter we covered:

\begin{itemize}[nosep]
  \item Why single-cell resolution is necessary and what bulk RNA-seq misses.
  \item The historical development of scRNA-seq technologies from 2009 to the
    present.
  \item Cell isolation methods: FACS, microfluidics, droplet-based (10x
    Chromium, Drop-seq, inDrop), combinatorial indexing (sci-RNA-seq, SPLiT-seq),
    and plate-based (Smart-seq2/3).
  \item The 10x Genomics Chromium system in detail: GEM formation, cell barcodes,
    UMIs, 3$'$ vs.\ 5$'$ chemistry, Cell Ranger, and the Flex assay.
  \item Computational challenges: sparsity/dropout, doublet detection, and
    ambient RNA contamination.
  \item The complete analysis workflow: raw data processing (Cell Ranger, STARsolo,
    alevin-fry), QC filtering, normalization (scran, SCTransform), feature
    selection (HVGs), dimensionality reduction (PCA, t-SNE, UMAP), graph-based
    clustering (Leiden/Louvain), cell type annotation, differential expression
    (pseudobulk recommended), trajectory analysis (Monocle3, scVelo),
    cell--cell communication (CellChat, NicheNet), and gene regulatory networks
    (SCENIC).
  \item Scanpy vs.\ Seurat frameworks, and dataset integration methods.
  \item Large-scale cell atlases: Human Cell Atlas, Tabula Sapiens, Tabula Muris.
  \item Practical workflows in both Snakemake and Nextflow.
\end{itemize}

The next chapter extends our view from dissociated single cells to cells in
their native tissue context with spatial transcriptomics (\cref{ch:spatial}).
